% ═══════════════════════════════════════════════════════════════════════════
%  فصل ۱ — درآمد، پرسش‌ها و چارچوب مفهومی
% ═══════════════════════════════════════════════════════════════════════════
\chapter{درآمد، پرسش‌ها و چارچوب مفهومی}
\label{ch:introduction}

\begin{tcolorbox}[mybox, title={چکیده‌ی فصل}]
این فصل به طرح مسئله‌ی اصلی پژوهش، تعریف مفاهیم کلیدی، مرور ادبیات نظری، و بیان چارچوب مفهومی کتاب اختصاص دارد. پرسش محوری آن است که پدیده‌ی «نجات‌طلبی» — یعنی دعوت از نیروی خارجی برای رهایی از ستم داخلی — چه ماهیتی دارد، چه الگوهایی در تاریخ بروز داده، و با چه معیارهایی می‌توان آن را ارزیابی کرد.
\end{tcolorbox}

% ─────────────────────────────────────────────────────────────────────────
\section{طرح مسئله}
\label{sec:ch01-problem}
% ─────────────────────────────────────────────────────────────────────────

در سراسر تاریخ بشر، گروه‌های انسانی در مواجهه با ستم، سرکوب یا بحران‌های وجودی، راه‌حل‌های گوناگونی جسته‌اند. یکی از مهم‌ترین و بحث‌برانگیزترین این راه‌حل‌ها، \concept{دعوت از نیروی خارجی} برای مداخله — اعم از نظامی، سیاسی یا اقتصادی — بوده است.

این پدیده در زبان این پژوهش \concept{نجات‌طلبی} نامیده می‌شود: فرایندی که طی آن بخشی از یک جامعه‌ی سیاسی (ملت، قوم، طبقه، فرقه‌ی مذهبی یا نخبگان سیاسی) از قدرتی بیرونی می‌خواهد تا در امور داخلی آنان مداخله کرده و وضعیت موجود را تغییر دهد.%
\footnote{این تعریف وسیع‌تر از مفهوم \lr{Intervention by Invitation} در حقوق بین‌الملل است. رجوع کنید به: \lr{Doswald-Beck, Louise. "The Legal Validity of Military Intervention by Invitation of the Government." \textit{British Year Book of International Law} 56, no. 1 (1985): 189--252.}}

\thinker{هدایت}، در «بوف کور»، تصویری نمادین از انزوای ملتی ترسیم می‌کند که راه نجاتی نمی‌یابد. اما در واقعیت تاریخی، ملت‌ها همیشه «بوف کور» نمانده‌اند — گاه دست به دامان دیگری شده‌اند. پرسش اصلی ما این است:

\begin{tcolorbox}[defbox, title={پرسش‌های محوری پژوهش}]
\begin{enumerate}[nosep, label=\textbf{پ\arabic*}.]
  \item آیا نجات‌طلبی یک الگوی تکرارشونده در تاریخ است و اگر هست، چه عوامل ساختاری آن را ایجاد می‌کند؟
  \item چه تمایز معناداری میان «دعوت مشروع برای مداخله‌ی انسان‌دوستانه» و «خیانت به حاکمیت ملی» وجود دارد؟
  \item الگوی نتایج مداخلات خارجی — اعم از فوری و بلندمدت — چیست و آیا الگوی قابل تعمیمی وجود دارد؟
  \item نقش عوامل روان‌شناسی جمعی (ناامیدی، آرزواندیشی، هویت‌طلبی) در شکل‌گیری نجات‌طلبی چیست؟
  \item آیا فعالیت‌های میسیونری، ایدئولوژیک و فرهنگی در تسهیل پذیرش مداخله‌ی خارجی نقش دارند؟
\end{enumerate}
\end{tcolorbox}

% ─────────────────────────────────────────────────────────────────────────
\section{تعریف مفاهیم کلیدی}
\label{sec:ch01-definitions}
% ─────────────────────────────────────────────────────────────────────────

\subsection{نجات‌طلبی}
\label{subsec:ch01-def-salvation}

\concept{نجات‌طلبی} (\lr{Salvation-Seeking}) در این پژوهش به معنای هرگونه تلاش سازمان‌یافته یا نیمه‌سازمان‌یافته‌ی بخشی از یک واحد سیاسی برای جلب مداخله‌ی قدرت خارجی به‌منظور تغییر وضع موجود سیاسی، نظامی یا اجتماعی تعریف می‌شود.

این تعریف طیف وسیعی را شامل می‌شود:

\begin{longtable}{
  >{\raggedright\arraybackslash}p{3cm}
  >{\raggedright\arraybackslash}p{4.5cm}
  >{\raggedright\arraybackslash}p{6cm}}
\caption{گونه‌شناسی نجات‌طلبی}\label{tab:ch01-typology}\\
\toprule
\rowcolor{PrimaryDeep}
\textcolor{white}{\textbf{نوع}} &
\textcolor{white}{\textbf{ماهیت درخواست}} &
\textcolor{white}{\textbf{نمونه‌ی تاریخی}} \\
\midrule
\endfirsthead

\toprule
\rowcolor{PrimaryDeep}
\textcolor{white}{\textbf{نوع}} &
\textcolor{white}{\textbf{ماهیت درخواست}} &
\textcolor{white}{\textbf{نمونه‌ی تاریخی}} \\
\midrule
\endhead

\midrule
\multicolumn{3}{l}{\small\itshape ادامه در صفحه‌ی بعد...}\\
\endfoot

\bottomrule
\endlastfoot

نظامی مستقیم &
درخواست لشکرکشی و اشغال &
دعوت پارلمان انگلستان از ویلیام اورانژ (۱۶۸۸)%
\footnote{\lr{Israel, Jonathan. \textit{The Anglo-Dutch Moment}. Cambridge UP, 1991.}} \\
\rowcolor{NeutralLight}
نظامی غیرمستقیم &
درخواست تسلیحات و آموزش &
کمک‌های امریکا به مجاهدین افغان (۱۹۷۹–۱۹۸۹)%
\footnote{\lr{Coll, Steve. \textit{Ghost Wars}. Penguin, 2004.}} \\
دیپلماتیک-سیاسی &
درخواست فشار سیاسی و تحریم &
درخواست اپوزیسیون آفریقای جنوبی برای تحریم آپارتاید%
\footnote{\lr{Klotz, Audie. \textit{Norms in International Relations}. Cornell UP, 1995.}} \\
\rowcolor{NeutralLight}
انسان‌دوستانه &
درخواست حمایت سازمان‌های بین‌المللی &
ایزدی‌های عراق (۲۰۱۴)%
\footnote{\lr{UNHCR. "Iraq Emergency." 2014.}} \\
فرهنگی-ایدئولوژیک &
پذیرش و ترویج ایدئولوژی خارجی &
نفوذ پان‌اسلاویسم در بالکان%
\footnote{\lr{Jelavich, Barbara. \textit{History of the Balkans}. Cambridge UP, 1983.}} \\
\rowcolor{NeutralLight}
اقتصادی &
واگذاری امتیازات اقتصادی در قبال حمایت &
امتیازنامه‌ی رویتر (۱۸۷۲)%
\footnote{آدمیت، فریدون. \textit{اندیشه‌ی ترقی و حکومت قانون: عصر سپهسالار}. خوارزمی، ۱۳۵۱.} \\

\end{longtable}

\subsection{مداخله‌ی خارجی}
\label{subsec:ch01-def-intervention}

\concept{مداخله‌ی خارجی} (\lr{Foreign Intervention}) در ادبیات روابط بین‌الملل به «هرگونه اقدام اجباری یا نیمه‌اجباری یک دولت یا سازمان بین‌المللی در امور داخلی دولتی دیگر» اشاره دارد.%
\footnote{\lr{Vincent, R.J. \textit{Nonintervention and International Order}. Princeton UP, 1974, p.~13.}}

\thinker{هدلی بول} (\lr{Hedley Bull}) مداخله را «اعمال نفوذ اجباری» تعریف می‌کند که «حق خودمختاری واحد سیاسی هدف را نقض می‌کند.»%
\footnote{\lr{Bull, Hedley. "Introduction." In \textit{Intervention in World Politics}, ed. Bull. Oxford UP, 1984, p.~1.}}

تمایز مهمی که باید روشن شود، تمایز میان \concept{مداخله‌ی تقاضامحور} (\lr{Intervention by Invitation}) و \concept{مداخله‌ی یک‌جانبه} (\lr{Unilateral Intervention}) است. در حالت نخست، حداقل بخشی از جامعه‌ی هدف خواهان مداخله است؛ در حالت دوم، مداخله‌گر بدون دعوت وارد می‌شود.

\subsection{موفقیت و شکست: مسئله‌ی تعریف}
\label{subsec:ch01-def-success}

\begin{tcolorbox}[defbox, title={معضل تعریف «موفقیت»}]
«موفقیت» مداخله‌ی خارجی مفهومی ذاتاً چندبُعدی و وابسته به منظر (\lr{perspective-dependent}) است. آنچه از دید مداخله‌کننده موفقیت محسوب شود، ممکن است از دید ملت هدف شکست باشد و بالعکس.%
\footnote{\lr{Regan, Patrick M. \textit{Civil Wars and Foreign Powers}. U of Michigan P, 2000, p.~45.}}
\end{tcolorbox}

برای پرهیز از ابهام، در این پژوهش \concept{موفقیت} در سه بُعد سنجیده می‌شود:

\begin{enumerate}[nosep, label=\textbf{\arabic*.}]
  \item \textbf{بُعد فوری (۰ تا ۵ سال):} آیا هدف اعلام‌شده‌ی دعوت‌کنندگان محقق شد؟ (سقوط رژیم، خاتمه‌ی ستم، استقلال و غیره)
  \item \textbf{بُعد میان‌مدت (۵ تا ۲۵ سال):} آیا نظم جدید پایدار ماند؟ آیا وابستگی به مداخله‌گر تشدید شد؟
  \item \textbf{بُعد بلندمدت (بیش از ۲۵ سال):} آیا جامعه به خودمختاری و خوداتکایی رسید؟ آیا استقلال و حاکمیت ملی حفظ شد؟
\end{enumerate}

% ─────────────────────────────────────────────────────────────────────────
\section{مرور ادبیات نظری}
\label{sec:ch01-literature}
% ─────────────────────────────────────────────────────────────────────────

ادبیات نظری مرتبط با موضوع این پژوهش در چهار حوزه‌ی اصلی قابل دسته‌بندی است:

\subsection{نظریه‌های روابط بین‌الملل}
\label{subsec:ch01-lit-ir}

\subsubsection{واقع‌گرایی}
مکتب واقع‌گرایی (\lr{Realism}) مداخله‌ی خارجی را ابزاری در خدمت منافع ملی مداخله‌کننده تلقی می‌کند. از نگاه \thinker{هانس مورگنتاو} (\lr{Morgenthau})، «هیچ دولتی برای نجات ملت دیگر وارد جنگ نمی‌شود مگر آنکه منافعش ایجاب کند.»%
\footnote{\lr{Morgenthau, Hans J. \textit{Politics Among Nations}. 5th ed. Knopf, 1978, p.~10.}}

\thinker{کنت والتز} (\lr{Kenneth Waltz}) نیز با تأکید بر ساختار آنارشیک نظام بین‌الملل استدلال می‌کند که دولت‌ها ذاتاً به‌دنبال افزایش قدرت نسبی خود هستند و «نجات ملت‌ها» اغلب پوشش ایدئولوژیک برای توسعه‌طلبی است.%
\footnote{\lr{Waltz, Kenneth N. \textit{Theory of International Politics}. Waveland, 1979, pp.~102--128.}}

\subsubsection{لیبرالیسم}
در مقابل، مکتب لیبرال (\lr{Liberalism}) امکان مداخله‌ی انسان‌دوستانه‌ی واقعی را می‌پذیرد. \thinker{مایکل دویل} (\lr{Michael Doyle}) با تکیه بر سنت کانتی، استدلال می‌کند که دموکراسی‌ها می‌توانند (و گاه باید) برای دفاع از حقوق بشر مداخله کنند.%
\footnote{\lr{Doyle, Michael W. "Kant, Liberal Legacies, and Foreign Affairs." \textit{Philosophy and Public Affairs} 12, no.~3 (1983): 205--235.}}

\subsubsection{سازه‌انگاری}
\thinker{مارتا فینه‌مور} (\lr{Martha Finnemore}) نشان می‌دهد که \concept{هنجارهای بین‌المللی} درباره‌ی مداخله در طول زمان تحول یافته‌اند: از حمایت از مسیحیان تحت ستم در سده‌ی نوزدهم تا مداخله‌ی انسان‌دوستانه‌ی فراگیر در سده‌ی بیستم.%
\footnote{\lr{Finnemore, Martha. \textit{The Purpose of Intervention}. Cornell UP, 2003.}}

\subsection{نظریه‌های پسااستعماری}
\label{subsec:ch01-lit-postcolonial}

\thinker{ادوارد سعید} (\lr{Edward Said}) در اثر بنیادین «شرق‌شناسی» نشان می‌دهد که چگونه غرب از طریق بازنمایی «شرقِ» ناتوان و نیازمند نجات، مداخلات خود را مشروعیت بخشیده است. از این منظر، حتی «دعوت» بومیان نیز ممکن است محصول هژمونی فرهنگی مداخله‌گر باشد.%
\footnote{\lr{Said, Edward W. \textit{Orientalism}. Vintage, 1978.}}

\thinker{فرانتس فانون} (\lr{Frantz Fanon}) نیز در «پوست سیاه، ماسک‌های سفید» و «دوزخیان روی زمین» بر پدیده‌ی \concept{درونی‌سازی سلطه} تأکید می‌کند: جایی که ملت تحت ستم، خود را ناتوان از رهایی درونی می‌بیند و «نجات‌دهنده» را در بیرون جست‌وجو می‌کند.%
\footnote{\lr{Fanon, Frantz. \textit{The Wretched of the Earth}. Trans. Constance Farrington. Grove, 1963.}}

\subsection{ادبیات فارسی}
\label{subsec:ch01-lit-persian}

در ادبیات فارسی، موضوع مداخله‌ی خارجی عمدتاً ذیل مفاهیم «استعمار»، «امپریالیسم» و «وابستگی» بررسی شده است. آثار کلیدی عبارت‌اند از:

\begin{itemize}[nosep, rightmargin=1em]
  \item \textbf{احمد اشرف} و \textbf{علی بنوعزیزی}،  «طبقات اجتماعی، دولت و انقلاب در ایران» (۱۳۸۶): تحلیل طبقاتی از مداخلات خارجی در ایران.%
  \footnote{اشرف، احمد و بنوعزیزی، علی. \textit{طبقات اجتماعی، دولت و انقلاب در ایران}. نیلوفر، ۱۳۸۶.}
  \item \textbf{یرواند آبراهامیان}، «ایران بین دو انقلاب» (۱۹۸۲): بررسی نقش قدرت‌های خارجی در تحولات سیاسی ایران.%
  \footnote{\lr{Abrahamian, Ervand. \textit{Iran Between Two Revolutions}. Princeton UP, 1982.}}
  \item \textbf{همایون کاتوزیان}، «اقتصاد سیاسی ایران» (۱۹۸۱): تحلیل اقتصاد سیاسی وابستگی.%
  \footnote{\lr{Katouzian, Homa. \textit{The Political Economy of Modern Iran}. Macmillan, 1981.}}
  \item \textbf{فریدون آدمیت}، مجموعه‌ی آثار درباره‌ی مشروطه و اندیشه‌ی ترقی.
\end{itemize}

% ─────────────────────────────────────────────────────────────────────────
\section{چارچوب مفهومی پژوهش}
\label{sec:ch01-framework}
% ─────────────────────────────────────────────────────────────────────────

بر مبنای مرور ادبیات، چارچوب مفهومی زیر برای تحلیل پدیده‌ی نجات‌طلبی پیشنهاد می‌شود:

\begin{figure}[htbp]
\centering
\begin{tikzpicture}[
  node distance=1.5cm and 2cm,
  every node/.style={font=\small},
  box/.style={
    draw=PrimaryMid, fill=BgCool, rounded corners=4pt,
    text width=3.2cm, align=center, minimum height=1.2cm,
    font=\small\bfseries
  },
  arrow/.style={->, >=Stealth, thick, PrimaryMid},
  label style/.style={font=\scriptsize, text=NeutralDark}
]

% ── ردیف بالا: عوامل ──
\node[box, fill=BgWarm] (struct)
  {\rl{عوامل ساختاری}\\%
   \rl{\scriptsize استبداد، ستم، بحران}};

\node[box, fill=BgWarm, left=2cm of struct] (psych)
  {\rl{عوامل روان‌شناختی}\\%
   \rl{\scriptsize ناامیدی، آرزواندیشی}};

\node[box, fill=BgWarm, right=2cm of struct] (ext)
  {\rl{عوامل خارجی}\\%
   \rl{\scriptsize فرصت‌طلبی قدرت‌ها}};

% ── ردیف میانی: پدیده ──
\node[box, fill=AccentGold!20, below=2cm of struct,
      text width=4cm, minimum height=1.5cm] (salv)
  {\rl{\large نجات‌طلبی}};

% ── ردیف پایین: نتایج ──
\node[box, fill=BgSuccess, below left=2.5cm and 1cm of salv] (suc)
  {\rl{موفقیت}\\%
   \rl{\scriptsize رهایی، استقلال}};

\node[box, fill=BgFail, below right=2.5cm and 1cm of salv] (fail)
  {\rl{شکست}\\%
   \rl{\scriptsize وابستگی، خسران}};

\node[box, fill=AccentAmber!20, below=2.5cm of salv] (mix)
  {\rl{نتایج مختلط}\\%
   \rl{\scriptsize آزادی + وابستگی}};

% ── فلش‌ها ──
\draw[arrow] (struct) -- (salv);
\draw[arrow] (psych) -- (salv);
\draw[arrow] (ext) -- (salv);
\draw[arrow] (salv) -- (suc);
\draw[arrow] (salv) -- (fail);
\draw[arrow] (salv) -- (mix);

% ── برچسب‌های فلش ──
\node[label style, above right=0.1cm and 0.1cm of salv.south west]
  {\rl{بُعد فوری}};
\node[label style, above left=0.1cm and 0.1cm of salv.south east]
  {\rl{بُعد بلندمدت}};

\end{tikzpicture}
\caption{چارچوب مفهومی تحلیل نجات‌طلبی}
\label{fig:ch01-framework}
\end{figure}

این چارچوب بر سه دسته‌ی عوامل تأکید دارد:

\begin{enumerate}[nosep, label=\textbf{\alph*.}]
  \item \textbf{عوامل ساختاری-سیاسی:} ماهیت رژیم حاکم، شدت سرکوب، بحران مشروعیت، و فروپاشی نهادهای داخلی.
  \item \textbf{عوامل روان‌شناسی جمعی:} ناامیدی از تغییر درونی، آرزواندیشی درباره‌ی «نجات‌دهنده»، و تروماهای تاریخی.
  \item \textbf{عوامل بین‌المللی:} منافع ژئوپولیتیک قدرت‌های خارجی، وجود یا فقدان هنجارهای بین‌المللی حمایتی، و ساختار نظام بین‌الملل.
\end{enumerate}

% ─────────────────────────────────────────────────────────────────────────
\section{گونه‌شناسی مداخلات از منظر عامل دعوت‌کننده}
\label{sec:ch01-typology-actor}
% ─────────────────────────────────────────────────────────────────────────

یکی از پرسش‌های مهم آن است که \textit{چه کسی} دعوت‌کننده است. این تمایز اهمیت حقوقی و اخلاقی فراوانی دارد:

\begin{tcolorbox}[casebox, title={جدول تحلیلی: عامل دعوت‌کننده}]
\begin{adjustbox}{max width=\linewidth}
\begin{tabular}{
  >{\raggedleft\arraybackslash}p{2.5cm}
  >{\raggedleft\arraybackslash}p{3.5cm}
  >{\raggedleft\arraybackslash}p{3cm}
  >{\raggedleft\arraybackslash}p{3.5cm}}
\toprule
\rowcolor{PrimaryMid}
\textcolor{white}{\textbf{عامل}} &
\textcolor{white}{\textbf{مشروعیت حقوقی}} &
\textcolor{white}{\textbf{مشروعیت اخلاقی}} &
\textcolor{white}{\textbf{نمونه}} \\
\midrule
حکومت مستقر &
بالا (اصل حاکمیت) &
بحث‌برانگیز &
دعوت اسد از روسیه (۲۰۱۵) \\
\rowcolor{NeutralLight}
اپوزیسیون سیاسی &
پایین تا متوسط &
بحث‌برانگیز &
شورای ملی لیبی (۲۰۱۱) \\
نخبگان نظامی &
بسیار پایین &
بسیار بحث‌برانگیز &
ژنرال‌های عراق (۲۰۰۳) \\
\rowcolor{NeutralLight}
اقلیت قومی-مذهبی &
متغیر (حق تعیین سرنوشت) &
بالا (اگر ستم مستند باشد) &
کردهای عراق (۱۹۹۱) \\
مردم (جنبش عمومی) &
متغیر &
بالاترین &
کوزوو (۱۹۹۹)%
\footnote{\lr{Independent International Commission on Kosovo. \textit{The Kosovo Report}. Oxford UP, 2000.}} \\
\bottomrule
\end{tabular}
\end{adjustbox}
\end{tcolorbox}

% ─────────────────────────────────────────────────────────────────────────
\section{تمایز مفهومی: نجات‌طلبی در برابر مفاهیم مشابه}
\label{sec:ch01-distinctions}
% ─────────────────────────────────────────────────────────────────────────

نجات‌طلبی با چندین مفهوم مشابه اما متمایز در ارتباط است که تفکیک آن‌ها ضروری است:

\subsection{نجات‌طلبی در برابر خیانت}

در گفتمان ناسیونالیستی، هرگونه ارتباط با قدرت خارجی «خیانت» تلقی می‌شود. اما این خوانش ساده‌انگارانه است. \thinker{بندیکت اندرسون} (\lr{Benedict Anderson}) نشان داده که مفهوم «ملت» خود ساختاری تخیلی (\lr{imagined community}) است و بنابراین مرز میان «خودی» و «بیگانه» نیز ساختاری و تاریخی است، نه ذاتی.%
\footnote{\lr{Anderson, Benedict. \textit{Imagined Communities}. Verso, 1983.}}

\subsection{نجات‌طلبی در برابر همکاری نظامی متقابل}

پیمان‌های نظامی دوجانبه (مانند ناتو) مبتنی بر «برابری حقوقی» دولت‌هایند، در حالی که نجات‌طلبی اغلب رابطه‌ای نابرابر است که طرف داخلی از موضع ضعف عمل می‌کند.

\subsection{نجات‌طلبی در برابر امپریالیسم}

\begin{tcolorbox}[defbox, title={تمایز ظریف}]
امپریالیسم فرایندی «از بالا به پایین» است که قدرت بزرگ به‌دنبال سلطه است. نجات‌طلبی فرایندی «از پایین به بالا» است که بخشی از جامعه‌ی تحت ستم به‌دنبال نجات است. اما در عمل، این دو اغلب درهم‌تنیده‌اند: امپریالیسم از نجات‌طلبی بهره می‌برد و نجات‌طلبی ناخواسته راه را برای امپریالیسم هموار می‌کند.
\end{tcolorbox}

% ─────────────────────────────────────────────────────────────────────────
\section{بازه‌ی زمانی و جغرافیایی پژوهش}
\label{sec:ch01-scope}
% ─────────────────────────────────────────────────────────────────────────

\begin{figure}[htbp]
\centering
\begin{tikzpicture}[
  x=0.022\linewidth,
  every node/.style={font=\scriptsize}
]

% خط زمان
\draw[timeline arrow] (0,0) -- (60,0);

% دوره‌ها
\fill[EraAncient, opacity=0.3] (0,-0.5) rectangle (20,0.5);
\fill[EraMedieval, opacity=0.3] (15,-0.5) rectangle (40,0.5);
\fill[EraEarlyMod, opacity=0.3] (35,-0.5) rectangle (60,0.5);
\fill[EraModern, opacity=0.3] (45,-0.5) rectangle (80,0.5);
\fill[EraContemp, opacity=0.3] (60,-0.5) rectangle (100,0.5);

% برچسب‌ها
\node[era label, fill=EraAncient] at (10, 1.2) {\rl{باستان}};
\node[era label, fill=EraMedieval] at (25, 1.2) {\rl{میانه}};
\node[era label, fill=EraEarlyMod] at (40, 1.2) {\rl{نوین اولیه}};
\node[era label, fill=EraModern] at (50, 1.2) {\rl{مدرن}};
\node[era label, fill=EraContemp] at (55, 1.2) {\rl{معاصر}};

% رویدادهای کلیدی
\node[event node, below=1cm] at (5, -0.5)
  {\rl{حمله‌ی اسکندر}\\%
   \rl{۳۳۴--۳۲۳ ق.م}};

\node[event node, below=1cm] at (15, -0.5)
  {\rl{فتح عرب}\\%
   \rl{۶۳۶--۶۵۱ م.}};

\node[event node, below=1cm] at (25, -0.5)
  {\rl{حمله‌ی مغول}\\%
   \rl{۱۲۱۹--۱۲۵۸ م.}};

\node[event node, below=1cm] at (35, -0.5)
  {\rl{انقلاب شکوهمند}\\%
   \rl{۱۶۸۸ م.}};

\node[event node, below=1cm] at (45, -0.5)
  {\rl{جنگ‌های قاجار-روس}\\%
   \rl{۱۸۰۴--۱۸۲۸ م.}};

\node[event node, below=1cm] at (55, -0.5)
  {\rl{مداخلات ۲۰۰۰–۲۰۲۰}\\%
   \rl{عراق، لیبی، سوریه}};

\end{tikzpicture}
\caption{خط زمان بازه‌ی پژوهش}
\label{fig:ch01-timeline}
\end{figure}

% ─────────────────────────────────────────────────────────────────────────
\section{ماتریس تحلیلی: ابعاد ارزیابی مداخلات}
\label{sec:ch01-matrix}
% ─────────────────────────────────────────────────────────────────────────

برای ارزیابی نظام‌مند هر مورد تاریخی، ماتریس زیر به‌کار گرفته می‌شود:

\begin{table}[htbp]
\centering
\caption{ماتریس ارزیابی چندبُعدی مداخلات خارجی}\label{tab:ch01-matrix}
\begin{adjustbox}{max width=\linewidth}
\begin{tabular}{
  >{\raggedleft\arraybackslash\bfseries}p{2.5cm}
  >{\raggedleft\arraybackslash}p{3cm}
  >{\raggedleft\arraybackslash}p{3cm}
  >{\raggedleft\arraybackslash}p{3cm}
  >{\raggedleft\arraybackslash}p{3cm}}
\toprule
\rowcolor{PrimaryDeep}
\textcolor{white}{بُعد ارزیابی} &
\textcolor{white}{شاخص فوری (۵>)} &
\textcolor{white}{شاخص میان‌مدت} &
\textcolor{white}{شاخص بلندمدت} &
\textcolor{white}{منبع سنجش} \\
\midrule
نظامی-امنیتی &
پایان خشونت &
ثبات امنیتی &
خودکفایی دفاعی &
آمار تلفات، شاخص صلح \\
\rowcolor{NeutralLight}
سیاسی-حکمرانی &
تغییر رژیم &
نهادسازی &
حکمرانی خوب &
\lr{Freedom House, Polity V} \\
اقتصادی &
بازسازی فوری &
رشد اقتصادی &
استقلال اقتصادی &
\lr{GDP, HDI} \\
\rowcolor{NeutralLight}
اجتماعی-فرهنگی &
کاهش تنش قومی &
انسجام اجتماعی &
هویت ملی پایدار &
پیمایش‌ها، شاخص گینی \\
حاکمیت ملی &
حفظ تمامیت ارضی &
استقلال سیاست خارجی &
عدم وابستگی ساختاری &
معاهدات، پایگاه‌های نظامی \\
\bottomrule
\end{tabular}
\end{adjustbox}
\end{table}

% ─────────────────────────────────────────────────────────────────────────
\section{نقشه‌ی مفهومی کتاب}
\label{sec:ch01-roadmap}
% ─────────────────────────────────────────────────────────────────────────

\begin{figure}[htbp]
\centering
\begin{tikzpicture}[
  mindmap,
  every node/.style={concept, font=\small\bfseries},
  root concept/.append style={
    concept color=PrimaryDeep, text=white,
    minimum size=3.5cm, font=\normalsize\bfseries
  },
  level 1 concept/.append style={
    concept color=PrimaryMid, text=white,
    minimum size=2.2cm, font=\small\bfseries
  },
  level 2 concept/.append style={
    concept color=PrimaryLight!60, text=PrimaryDeep,
    minimum size=1.5cm, font=\scriptsize
  },
  level 1/.append style={sibling angle=72, level distance=4.5cm},
  level 2/.append style={sibling angle=45, level distance=3cm},
]

\node[root concept] {\rl{نجات‌طلبی}}
  child[concept color=EraAncient!80] {
    node[concept] {\rl{مبانی نظری}}
    child { node[concept] {\rl{حقوق بین‌الملل}} }
    child { node[concept] {\rl{الگوهای تحلیلی}} }
  }
  child[concept color=AccentTeal!80] {
    node[concept] {\rl{ایران}}
    child { node[concept] {\rl{باستان}} }
    child { node[concept] {\rl{معاصر}} }
  }
  child[concept color=EraEarlyMod!80] {
    node[concept] {\rl{اروپا}}
    child { node[concept] {\rl{انقلاب شکوهمند}} }
    child { node[concept] {\rl{استعمار}} }
  }
  child[concept color=EraContemp!80] {
    node[concept] {\rl{معاصر جهانی}}
    child { node[concept] {\rl{لیبی، سوریه}} }
    child { node[concept] {\rl{اویغورها، ایزدی‌ها}} }
  }
  child[concept color=AccentPurple!80] {
    node[concept] {\rl{تحلیل‌ها}}
    child { node[concept] {\rl{روان‌شناسی}} }
    child { node[concept] {\rl{اقتصاد}} }
  };

\end{tikzpicture}
\caption{نقشه‌ی ذهنی ساختار کتاب}
\label{fig:ch01-mindmap}
\end{figure}

% ─────────────────────────────────────────────────────────────────────────
\section{جمع‌بندی فصل}
\label{sec:ch01-summary}
% ─────────────────────────────────────────────────────────────────────────

در این فصل، مسئله‌ی محوری پژوهش — یعنی پدیده‌ی \concept{نجات‌طلبی} — تعریف و مرزبندی شد. مرور ادبیات نشان داد که این پدیده در تقاطع چندین رشته‌ی علمی قرار دارد و هیچ رشته‌ی واحدی به‌تنهایی توانایی تحلیل جامع آن را ندارد. چارچوب مفهومی سه‌لایه‌ای (عوامل ساختاری، روان‌شناختی و بین‌المللی) و ماتریس ارزیابی پنج‌بُعدی (نظامی، سیاسی، اقتصادی، اجتماعی و حاکمیتی) به‌عنوان ابزارهای تحلیلی کتاب معرفی شدند. در فصل بعد، مبانی حقوقی مداخله‌ی خارجی در حقوق بین‌الملل بررسی خواهد شد.

\cleardoublepage